\chapter{Travaux futurs}\label{ch:futurs}

\begin{origine}
\textbf{Au départ}, les fragilités du \cref{ch:discussion} et les jalons du \cref{sec:decision}.
\textbf{Ce chapitre établit} l'ordre du travail qui vient, chaque pas avec la mesure qui le juge ; les questions de recherche ouvertes, chacune chiffrée par l'équation du jeton ; le seul geste qui revienne à Stone ; et les téléchargements qui pourraient servir, aucun n'étant nécessaire au premier programme.
\textbf{Il reste} à le faire.
\end{origine}

\section{Les mesures qui décident, d'abord}\label{sec:mesures-decident}

\begin{enumerate}
\item \textbf{La courbe $h(c)$ et la réutilisation $\rho$ de gpt-oss-120b}, sur ses propres traces de routage. Elles remplacent l'\cref{hyp:beta} dans tous les plafonds ; la mesure du 20b (\cref{fig:localite-mesuree}) dit déjà que la forme supposée est fausse. Témoin : la réplique du cache de llama.cpp rejouée sur les traces, qui doit redonner les nombres de son journal.
\item \textbf{La base llama.cpp de gpt-oss-120b en séries ABBA}, sur secteur, selon le protocole du \cref{sec:protocole-jugement} : c'est la première mesure qui peut tuer A (le rapport $\times 2{,}5$ à $\times 17$ de l'étude 07 repose sur une base estimée).
\item \textbf{Le banc Q0} : la qualité de B sans moteur, par fichiers MXFP4 statiques (base aimantée sur des valeurs E2M1 existantes) jugés par llama.cpp, et la frontière des requantifications uniformes. C'est la mesure qui décide de J2.
\item \textbf{L'écart entre deux moteurs justes sur gpt-oss}, qui recale le budget de qualité $\kappa$ du \cref{def:probleme}.
\end{enumerate}

\section{Le jalon J1 : construire A}\label{sec:construire-j1}

L'arbitre découpe le travail en cinq lots aux fichiers disjoints, chacun fini par un témoin, sa cassure et un seuil déclaré (\fichier{decisions/ARCHITECTURE.md}, § 6) :
\begin{description}
\item[F, format et codeur] : réécriture du GGUF d'origine en fichier \texttt{.trc} v1.0 ; témoin : chaque tronçon décodé égale l'original, octet pour octet.
\item[E, chemin des octets] : lecteur IOCP à la taille d'un expert vers la mémoire épinglée, DMA ; témoin : le débit de bout en bout ne dépasse aucun étage, et les octets arrivent exacts.
\item[K, noyaux] : décodeur rANS à table par voie, produit MXFP4, passage avant de gpt-oss ; témoin : les noyaux WGSL déjà mesurés, sur les mêmes octets (second chemin).
\item[R, référence et juge] : une référence processeur en double précision, le juge de qualité, le banc Q0.
\item[O, cache, ordonnanceur, moteur] : cache d'experts en VRAM, file des défauts, modes A0 et A ; témoin : logits identiques au bit entre le chemin codé et le chemin brut.
\end{description}
Le jalon se juge sur gpt-oss-120b : jetons/s, RAM et VRAM de pointe, octets lus selon le compteur du système, et rendement physique, c'est-à-dire le débit mesuré rapporté au plafond de l'\cref{eq:jeton} recalculé avec le succès compté par le moteur. Puis J1+ ajoute le préchargement d'une couche d'avance et un cache SLRU, et se juge contre la meilleure base llama.cpp mesurée.

\section{Les questions de recherche}\label{sec:questions-recherche}

\begin{table}[ht]
\centering\small
\caption[Les questions de recherche ouvertes]{Les questions de recherche ouvertes, et ce que chacune vaut d'après l'équation du jeton \calc~(\fichier{fondations/07-ce-qui-manque.md}, § 3.2, et \cref{sec:correction-b1}).}
\label{tab:recherche}
\begin{tabular}{@{}p{0.06\linewidth}p{0.36\linewidth}p{0.5\linewidth}@{}}
\toprule
& Question & Ce qu'elle vaut, et ce qui manque \\
\midrule
R1 & Le raffinement exact d'un format 4 bits natif, lié à l'étage (B) & en débit : B3, $\times 1{,}74$ au prix de 48,5\,\% d'usages servis par la base ; en latence : B1, plancher $+23\,\%$. Manque : la qualité (Q0) \\
R2 & L'allocation conjointe des bits et de l'étage & avec les débits de la machine, une matrice résidente mérite $\tfrac12 \log_2(209/2{,}58) = 3{,}2$ bits de plus qu'une matrice lue du disque ; aucun outil ne l'applique \\
R3 & Une politique de défaut sous budget de qualité & attendre, reporter, servir la base, re-router, défaut par défaut ; le gain vit entre le plancher et le plafond (A, LRU : 3,93 $\to$ 6,24, $\times 1{,}59$) \\
R4 & La spéculation sous réutilisation réelle & condition $\alpha > 1 - \rho$ (\cref{prop:spec-rho}) ; $\rho = 0{,}577$ sur gpt-oss-20b ; à mesurer sur le 120b, et pour $k > 1$ \\
R5 & Un code à 2 bits de deuxième génération décodable sans CUDA & décodeur de treillis écrit et mesuré à 124\,Go/s, codeur absent ; $B_e$ divisé par 1,9 (4,25 $\to$ 2,25 bits) ; qualité inconnue sur un modèle entraîné en 4 bits \\
R6 & Les arbres de vérification pour les hybrides à règle delta & 48\,Mio d'état par nœud gardé, 12\,Gio pour 256 nœuds ; bloque C sur les modèles Qwen \\
R7 & Un code à décodage en temps constant pour le régime résident & au mieux $+10\,\%$ de capacité de la VRAM, $+2\,\%$ de plafond ; basse priorité \\
\bottomrule
\end{tabular}
\end{table}

\section{Le geste qui revient à Stone}\label{sec:geste-stone}

Le plus gros levier de la machine n'est pas un logiciel (\cref{sec:bilan}). Un volume \texttt{D:} non chiffré, créé par un script déjà écrit et que Stone lancerait lui-même en administrateur (\fichier{outils/creer-disque-modeles.ps1}), porterait le plafond de A de 6,24 à \textbf{14,1 jetons/s} dans le scénario LRU calé si le disque y atteignait les 6,0\,Go/s de son annonce \calc~(\fichier{fondations/journaux/07-synthese/csv/architectures-gptoss120b.csv}, ligne \texttt{A\_sans\_perte\_disque\_non\_chiffre\_6}). C'est une borne haute : Intel RST VMD reste sur le chemin, et sa part dans les 53\,\% perdus n'est pas séparée. Au-delà, le lien reprend la main (58\,ms contre 67 pour le disque) : garder le lien en Gen 4 et la capacité de la VRAM redeviennent des leviers.

\section{Les téléchargements possibles}\label{sec:telechargements-futurs}

Aucun n'est nécessaire pour construire et juger A et B sur gpt-oss-120b, déjà sur la machine. Tout téléchargement de plus de 500\,Mo demande l'accord de Stone, avec le nom, la source, la taille et la licence (\cref{tab:telechargements-futurs}).

\begin{table}[ht]
\centering\small
\caption[Les téléchargements qui pourraient servir]{Les téléchargements qui pourraient servir, triés par nécessité \litt~(tailles lues par l'API Hugging Face ou NuGet, \fichier{fondations/07-ce-qui-manque.md}, § 5).}
\label{tab:telechargements-futurs}
\begin{tabular}{@{}rp{0.34\linewidth}rp{0.15\linewidth}p{0.27\linewidth}@{}}
\toprule
Rang & Nom & Taille & Licence & Pour \\
\midrule
1 & DirectStorage 1.4 (préversion) & 3,99\,Mo & Microsoft & comparer le chemin disque $\to$ carte \\
2 & Qwen3.6-35B-A3B UD-Q4\_K\_M & 22,13\,Go & Apache-2.0 & second MoE fin, trace et $\rho$ \\
3 & OLMoE-1B-7B Q4\_K\_M & $\approx 4{,}2$\,Go & Apache-2.0 & trace de routage rapide \\
4 & Llama-3.3-70B IQ3\_XXS + Llama-3.2-1B Q4\_K\_M & 28,28\,Go & Llama 3.3 / 3.2 & architecture C \\
5 & Qwen3.8-Flash-Next UD-Q2\_K\_XL & 78,87\,Go & qwen-community-1.0 & second modèle de B \\
6 & Qwen3.5-9B-MTP Q4\_K\_M & 5,87\,Go & Apache-2.0 & accord du MTP natif \\
\bottomrule
\end{tabular}
\end{table}

\section{Le mémoire lui-même}\label{sec:memoire-suite}

Deux chapitres s'ajouteront à la troisième partie, avant la discussion : \textbf{l'implémentation} (les modules de troncons, le format \texttt{.trc} au bit près, les noyaux et leurs témoins) et \textbf{les résultats} (chaque série du jalon J1 en figure, avec barres d'erreur, témoin et seuil déclaré). \barre{L'étude de l'entropie, que ce mémoire cite sans lui consacrer de chapitre parce qu'elle s'achevait chez un autre agent, recevra le sien} (elle l'a reçu : \cref{ch:entropie}), et la mesure de gpt-oss, une fois rendue, remplacera les figures provisoires du \cref{sec:premieres-mesures}.

\begin{faux}
\begin{itemize}
\item L'ordre proposé suppose que la base mesurée de llama.cpp ne dépasse pas le plafond de A ; si elle le dépassait, A serait à repenser avant d'être construit.
\item Les gains de la \cref{tab:recherche} sont des plafonds calculés sous l'\cref{hyp:beta} ; ils bougeront avec la courbe mesurée.
\item Le gain du volume non chiffré est une borne haute fondée sur une annonce de fabricant, pas sur une mesure.
\end{itemize}
\end{faux}
